% ---
% title: What investors hold
% tags: finance, inclusion-exclusion, venn-diagrams
% ---

\begin{problem}
	\begin{context}
		A \emph{stock} (share) is a small piece of ownership in a company. A
		\emph{bond} is a loan to a company or a government that pays interest. An
		\emph{investment fund unit} is a share in a fund that pools the money of
		many investors and invests it in many stocks or bonds at once.
	\end{context}
	Among the retail clients of a brokerage, a randomly chosen client holds
	stocks ($S$) with probability $0.45$, bonds ($B$) with probability $0.35$
	and investment fund units ($F$) with probability $0.30$. The client holds
	both stocks and bonds with probability $0.15$, stocks and fund units with
	probability $0.12$, bonds and fund units with probability $0.10$, and all
	three with probability $0.05$. Find the probability that the client holds
	(a) at least one of the three; (b) none of them; (c) exactly one of them;
	(d) stocks but no bonds.
\end{problem}

\begin{solution}
	\idea ``At least one'' is the union $S\cup B\cup F$, so we use
	inclusion--exclusion for three events. For ``exactly one'' we fill the
	Venn diagram with probabilities, from the inside out.

	\step{(a)} ``At least one of the three'' means the event $S\cup B\cup F$.
	For three events, inclusion--exclusion says
	\[ P(S\cup B\cup F)=P(S)+P(B)+P(F)-P(SB)-P(SF)-P(BF)+P(SBF). \]
	Why this form? Adding the three probabilities counts a client with two
	products twice, so we subtract the pairs. But then a client with all three
	products was added three times and subtracted three times, so we add
	$P(SBF)$ back once. With the numbers:
	\[ P(S\cup B\cup F)=0.45+0.35+0.30-0.15-0.12-0.10+0.05=0.78. \]

	\step{(b)} ``None'' is the complement of ``at least one'':
	\[ P\big((S\cup B\cup F)^c\big)=1-P(S\cup B\cup F)=1-0.78=0.22. \]

	\step{(c)} Here inclusion--exclusion alone is not enough: we need the
	probability of each region of the Venn diagram. We fill it from the inside
	out, because the inner regions are needed to compute the outer ones.

	\begin{minipage}{.54\linewidth}
		\begin{itemize}[nosep]
			\item center (all three): $P(SBF)=0.05$;
			\item $S$ and $B$ only (no $F$): $0.15-0.05=0.10$;
			\item $S$ and $F$ only: $0.12-0.05=0.07$;
			\item $B$ and $F$ only: $0.10-0.05=0.05$;
			\item only $S$: $0.45-0.10-0.07-0.05=0.23$;
			\item only $B$: $0.35-0.10-0.05-0.05=0.15$;
			\item only $F$: $0.30-0.07-0.05-0.05=0.13$;
			\item outside all circles: $0.22$ (part (b)).
		\end{itemize}
	\end{minipage}\hfill
	\begin{minipage}{.42\linewidth}
		\centering
		\begin{tikzpicture}[scale=.8]
			\draw[thick] (-2.2,-3.6) rectangle (4,2.1);
			\draw[navy,thick] (0,0) circle (1.5);
			\draw[teal,thick] (1.8,0) circle (1.5);
			\draw[gold,thick] (.9,-1.5) circle (1.5);
			\node at (-.7,.3) {\small 0.23};
			\node at (2.5,.3) {\small 0.15};
			\node at (.9,-2.3) {\small 0.13};
			\node at (.9,.6) {\scriptsize 0.10};
			\node at (-.02,-1.0) {\scriptsize 0.07};
			\node at (1.82,-1.0) {\scriptsize 0.05};
			\node at (.9,-.4) {\scriptsize 0.05};
			\node at (3.3,-3.1) {\small 0.22};
			\node[navy] at (-1.4,1.4) {$S$};
			\node[teal] at (3.2,1.4) {$B$};
			\node[gold] at (-.9,-2.7) {$F$};
			\node at (.9,1.8) {\small $\Om$};
		\end{tikzpicture}
	\end{minipage}

	\medskip
	Each pairwise probability, e.g.\ $P(SB)=0.15$, contains the center
	$0.05$ as well, so we subtract it to get the part with exactly those two
	products. Each circle, e.g.\ $S$, consists of four regions: only $S$, the
	two ``exactly two'' regions inside $S$, and the center. So only $S$ is what
	remains of $0.45$. The same in one formula:
	\[ P(\text{only }S)=P(S)-P(SB)-P(SF)+P(SBF)=0.45-0.15-0.12+0.05=0.23. \]
	``Exactly one'' consists of the three outer regions, which are mutually
	exclusive, so their probabilities add up:
	\[ P(\text{exactly one})=0.23+0.15+0.13=0.51. \]
	Check: all eight regions together give
	$0.23+0.15+0.13+0.10+0.07+0.05+0.05+0.22=1$.

	\step{(d)} ``Stocks but no bonds'' is the event $SB^c$. The event $S$ splits
	into two mutually exclusive parts: $SB$ (stocks and bonds) and $SB^c$
	(stocks, no bonds). Hence $P(S)=P(SB)+P(SB^c)$, and
	\[ P(SB^c)=P(S)-P(SB)=0.45-0.15=0.30. \]
	In the diagram: the two regions of $S$ outside $B$, $0.23+0.07=0.30$.
	\begin{answer}
		(a) $0.78$ \quad (b) $0.22$ \quad (c) $0.51$ \quad (d) $0.30$
	\end{answer}
	\begin{remark}
		\textbf{Common mistake:} computing ``only $S$'' as
		$0.45-0.15-0.12=0.18$. This subtracts the clients with all three
		products twice (they are in $SB$ and in $SF$), so $P(SBF)=0.05$ must be
		added back once. Filling the diagram from the center avoids this mistake.
	\end{remark}
\end{solution}
